\section{Fair Comparison Validation}
\label{sec:fair-comparison}

This section validates the experimental fairness of the three-system comparison on the 2WikiMultiHopQA benchmark: our Hybrid system, the LogicRAG baseline~\cite{logicrag2026}, and the LinearRAG baseline~\cite{linearrag2025}. An independent audit examined 37 edge cases across four categories: evaluation protocol, information access, statistical methodology, and selection bias. The audit found 23 fully handled, 12 partially handled, and 2 unhandled edge cases, with no issues affecting reported conclusions.

\subsection{Controlled Variables}
\label{sec:controlled-variables}

\paragraph{Identical question set.}
All three systems were evaluated on the same 1{,}000 questions (intersection = 1{,}000, zero exclusive IDs per system). Gold answers are identical across all systems, verified with zero mismatches. The 1{,}000-question subset derives from the LinearRAG HuggingFace dataset (\texttt{Zly0523/linear-rag}), itself a reformatting of 2WikiMultiHopQA's development set. Type distribution is near-proportional to the full 12{,}576-question set (maximum deviation: 1.6pp for \texttt{bridge\_comparison}).

\paragraph{Same LLM.}
All three systems use \texttt{llama-3.3-70b-versatile} via the Groq API with temperature 0. Both baselines were intentionally re-run with this model (replacing LogicRAG's original GPT-4o-mini and LinearRAG's original model) to isolate architectural differences from model quality variation. This deviation from original paper configurations is documented in reproduction notes.

\paragraph{Same judge model and prompt.}
All systems were judged by the same LLM-as-judge: \texttt{gemma3:12b} via Ollama, using the identical judge script (\texttt{judge\_hybrid\_eval.py}) with the same system prompt, evaluation prompt template, and word-level containment fast path. The containment precheck fires for 46--55\% of questions across systems (LogicRAG: 46.2\%, LinearRAG: 50.7\%, Hybrid: 54.5\%), reflecting system verbosity differences rather than judge bias.

\paragraph{Same metrics code.}
All comparison scripts import \texttt{compute\_em()}, \texttt{compute\_f1()}, and \texttt{compute\_containment()} from a single source (\texttt{evaluation/metrics.py}). The unified comparison script recomputes all metrics fresh from predicted and gold answers, ignoring any stale values stored in intermediate files.

\paragraph{Consistent empty-prediction handling.}
Empty predictions receive EM=0, F1=0, Containment=0, and \texttt{judge\_correct=False}. LogicRAG produces 127 ``Unknown''-like responses vs.\ Hybrid's 14 and LinearRAG's 10; these correctly receive low scores, reflecting genuine system behavior differences.

\subsection{Information Access Asymmetry}
\label{sec:info-access-asymmetry}

The three systems access fundamentally different information sources, which we disclose as a structural confound rather than a controlled variable.

\begin{table}[h]
\centering
\begin{tabular}{lllr}
\toprule
\textbf{System} & \textbf{Corpus} & \textbf{Access} & \textbf{LLM calls/Q} \\
\midrule
Hybrid & Live Wikipedia + DBpedia SPARQL & Open, live & 3--10 \\
LogicRAG & 6{,}119 benchmark snippets (254 char median) & Closed, fixed & 3--7 \\
LinearRAG & 658 full-article chunks (4{,}449 char median) & Closed, fixed & 1 \\
\bottomrule
\end{tabular}
\caption{Information access comparison. LLM calls/Q = typical number of LLM inference calls per question.}
\label{tab:info-access}
\end{table}

\paragraph{Favors Hybrid:}
\begin{itemize}
\item Unbounded live Wikipedia + DBpedia SPARQL access vs.\ fixed benchmark corpora
\item Higher computation budget: 3--10 LLM calls per question (routing, decomposition, sub-question solving, final synthesis) vs.\ LogicRAG's 3--7 calls and LinearRAG's single call
\item Hand-crafted routing rules encode domain knowledge (bridge patterns, comparison detection, keyword-to-property mappings) that baselines must learn implicitly
\item Temporal advantage: 2026 Wikipedia may contain information updated since the 2017 benchmark snapshot
\end{itemize}

\paragraph{Favors baselines:}
\begin{itemize}
\item Gold answers are guaranteed to exist in the benchmark corpus both baselines index
\item Surface-form alignment: benchmark answers match benchmark passage phrasing
\item No entity linking errors: both baselines use pre-resolved benchmark titles, while Hybrid uses DBpedia Spotlight (confidence $\geq 0.5$), which can fail or misclassify
\item LogicRAG's corpus is specifically designed for 2WikiMultiHopQA; LinearRAG's 658 full-article chunks provide richer per-entity context than LogicRAG's short snippets
\end{itemize}

\paragraph{Corpus architecture difference between baselines.}
LogicRAG indexes 6{,}119 short intro-paragraph snippets (median 254 characters, 2 sentences). LinearRAG indexes 658 concatenated full-article chunks (median 4{,}449 characters). While both use benchmark-provided text, LinearRAG gets substantially more context per entity. This architectural difference may partially explain LinearRAG's higher accuracy on compositional questions.

\subsection{Baseline Implementation Fidelity}
\label{sec:baseline-fidelity}

\paragraph{LogicRAG.}
We use the unmodified original LogicRAG codebase (\texttt{baselines/LogicRAG/src/models/logic\_rag.py}). The only change is the LLM backend: GPT-4o-mini $\to$ \texttt{llama-3.3-70b-versatile} via Groq. Embedding model (\texttt{all-MiniLM-L6-v2}), retrieval top-$k=3$, and maximum rounds=3 follow the code defaults. The warm-up analysis, dependency decomposition, topological sort, and iterative retrieval architecture are unchanged.

\paragraph{LinearRAG.}
We use the original LinearRAG codebase (\texttt{baselines/LinearRAG/src/LinearRAG.py}) with the same LLM swap. Embedding model (\texttt{all-mpnet-base-v2}), BFS iteration mode, and entity-graph retrieval architecture are unchanged. Predictions were generated by our re-run (not extracted from published results), confirmed by run logs and intermediate output files.

\subsection{Evaluation Design}
\label{sec:eval-design}

\paragraph{Dual-tier evaluation.}
We report results on two evaluation tiers:
\begin{enumerate}
\item \textbf{Hop-368}: 368 questions where all predicates in the shortest answer path exist in DBpedia (KG-answerable subset). This tier enables controlled comparison with SPARQL-only systems (Schema-Guided, GRASP).
\item \textbf{Full-1000}: All 1{,}000 questions. On the 632 non-KG-answerable questions, SPARQL-only systems have no text fallback, while Hybrid and LinearRAG do. Full-1000 numbers reflect overall system capability, not retrieval quality alone.
\end{enumerate}

\paragraph{Hop verification methodology.}
The 368-question subset was determined by live SPARQL queries against \texttt{https://dbpedia.org/sparql} (run: 2026-03-25), verifying that all predicates in each question's answer path exist and return values matching the ground truth. DBpedia ontology properties (e.g., \texttt{dbo:birthDate}, \texttt{dbo:director}) are generally stable, but re-verification on a different date may yield slightly different subset membership due to endpoint updates.

\paragraph{Answerability criterion consistency.}
Two answerability criteria exist: \texttt{is\_answerable} (375 questions, heuristic) and \texttt{hop\_verified\_answerable} (368 questions, ground truth). Their overlap is 366 questions (99.5\%). All comparison scripts consistently use the stricter \texttt{hop\_verified\_answerable} criterion.

\subsection{Statistical Methodology}
\label{sec:stat-methods}

\paragraph{McNemar test.}
Pairwise significance uses McNemar's chi-squared with Yates continuity correction: $\chi^2 = (|b - c| - 1)^2 / (b + c)$, where $b$ = System A correct \& B wrong, $c$ = A wrong \& B correct. Three independent implementations in the codebase were verified algebraically identical.

\paragraph{Multiple testing correction.}
Bonferroni correction adjusts $\alpha$ for the number of pairwise comparisons: $\alpha = 0.05/6 = 0.0083$ for hop-368 (6 pairs) and $\alpha = 0.05/3 = 0.0167$ for full-1000q (3 pairs).

\paragraph{Confidence intervals.}
Wilson score intervals for binary metrics (EM, Judge EM, Containment). Bootstrap percentile intervals ($n_{\text{boot}}=10{,}000$, seed=42) for continuous metrics (F1). These are standard non-parametric approaches.

\paragraph{Per-type analysis.}
Per-question-type breakdowns are descriptive, not inferential. Wilson CI widths range from 5.8pp (\texttt{compositional}, $n=413$ on 1000q) to 20.8pp (\texttt{inference}, $n=34$ on hop-368). Per-type differences between systems should be interpreted as exploratory patterns, not statistically confirmed effects, especially for subgroups with $n < 50$.

\subsection{Date Normalization}
\label{sec:date-norm}

The metrics code normalizes dates to canonical YYYY-MM-DD format before comparison, handling ISO (``1932-08-16''), DMY (``16 August 1932''), and MDY (``August 16, 1932'') formats. This follows the official 2WikiMultiHopQA evaluation protocol.

Hybrid produces 6 ISO-format date answers (from SPARQL returning raw DBpedia values) that benefit from this normalization, gaining approximately 3 additional EM matches (+0.3pp). LogicRAG and LinearRAG produce zero ISO-format dates. This asymmetry affects only raw EM; the primary Judge EM metric evaluates semantic equivalence and is unaffected.
